We present a multimodal information retrieval system for the CASTLE 2024 egocentric lifelogging dataset, comprising 416,542 keyframes sampled at 5-second intervals across four recording days and fifteen camera streams. Three retrieval approaches are implemented and evaluated. Approach A is a sparse BM25 baseline over ASR transcripts and \textbf{Florence-2} visual captions. Approach B performs dense visual retrieval with \textbf{SigLIP2} SO400M embeddings over the full frame corpus. Approach C, the primary contribution, combines sparse text, dense text, and dense visual signals through weighted late fusion, then applies cross-encoder reranking with \textbf{bge-reranker-v2-m3} over the top-50 candidates. Evaluation on 10 queries against 71 manually judged relevant frames yields mean Precision@10 scores of 0.070, 0.620, and 0.640 for Approaches A, B, and C respectively. The results show that dense visual retrieval is the strongest signal for this task, while fusion provides the best overall effectiveness. As an indexing contribution, we introduce an OCR hallucination gate that uses \textbf{SigLIP2} cosine similarity to remove 178,914 ungrounded OCR-bearing caption records before retrieval. An interactive Streamlit interface supports query selection, approach switching, and Rocchio relevance feedback for iterative refinement.
